% Options for packages loaded elsewhere
\PassOptionsToPackage{unicode}{hyperref}
\PassOptionsToPackage{hyphens}{url}
%
\documentclass[
  10pt,
]{article}
\usepackage{amsmath,amssymb}
\usepackage{iftex}
\ifPDFTeX
  \usepackage[T1]{fontenc}
  \usepackage[utf8]{inputenc}
  \usepackage{textcomp} % provide euro and other symbols
\else % if luatex or xetex
  \usepackage{unicode-math} % this also loads fontspec
  \defaultfontfeatures{Scale=MatchLowercase}
  \defaultfontfeatures[\rmfamily]{Ligatures=TeX,Scale=1}
\fi
\usepackage{lmodern}
\ifPDFTeX\else
  % xetex/luatex font selection
\fi
% Use upquote if available, for straight quotes in verbatim environments
\IfFileExists{upquote.sty}{\usepackage{upquote}}{}
\IfFileExists{microtype.sty}{% use microtype if available
  \usepackage[]{microtype}
  \UseMicrotypeSet[protrusion]{basicmath} % disable protrusion for tt fonts
}{}
\makeatletter
\@ifundefined{KOMAClassName}{% if non-KOMA class
  \IfFileExists{parskip.sty}{%
    \usepackage{parskip}
  }{% else
    \setlength{\parindent}{0pt}
    \setlength{\parskip}{6pt plus 2pt minus 1pt}}
}{% if KOMA class
  \KOMAoptions{parskip=half}}
\makeatother
\usepackage{xcolor}
\usepackage[margin=1in]{geometry}
\usepackage{color}
\usepackage{fancyvrb}
\newcommand{\VerbBar}{|}
\newcommand{\VERB}{\Verb[commandchars=\\\{\}]}
\DefineVerbatimEnvironment{Highlighting}{Verbatim}{commandchars=\\\{\}}
% Add ',fontsize=\small' for more characters per line
\newenvironment{Shaded}{}{}
\newcommand{\AlertTok}[1]{\textcolor[rgb]{1.00,0.00,0.00}{\textbf{#1}}}
\newcommand{\AnnotationTok}[1]{\textcolor[rgb]{0.38,0.63,0.69}{\textbf{\textit{#1}}}}
\newcommand{\AttributeTok}[1]{\textcolor[rgb]{0.49,0.56,0.16}{#1}}
\newcommand{\BaseNTok}[1]{\textcolor[rgb]{0.25,0.63,0.44}{#1}}
\newcommand{\BuiltInTok}[1]{\textcolor[rgb]{0.00,0.50,0.00}{#1}}
\newcommand{\CharTok}[1]{\textcolor[rgb]{0.25,0.44,0.63}{#1}}
\newcommand{\CommentTok}[1]{\textcolor[rgb]{0.38,0.63,0.69}{\textit{#1}}}
\newcommand{\CommentVarTok}[1]{\textcolor[rgb]{0.38,0.63,0.69}{\textbf{\textit{#1}}}}
\newcommand{\ConstantTok}[1]{\textcolor[rgb]{0.53,0.00,0.00}{#1}}
\newcommand{\ControlFlowTok}[1]{\textcolor[rgb]{0.00,0.44,0.13}{\textbf{#1}}}
\newcommand{\DataTypeTok}[1]{\textcolor[rgb]{0.56,0.13,0.00}{#1}}
\newcommand{\DecValTok}[1]{\textcolor[rgb]{0.25,0.63,0.44}{#1}}
\newcommand{\DocumentationTok}[1]{\textcolor[rgb]{0.73,0.13,0.13}{\textit{#1}}}
\newcommand{\ErrorTok}[1]{\textcolor[rgb]{1.00,0.00,0.00}{\textbf{#1}}}
\newcommand{\ExtensionTok}[1]{#1}
\newcommand{\FloatTok}[1]{\textcolor[rgb]{0.25,0.63,0.44}{#1}}
\newcommand{\FunctionTok}[1]{\textcolor[rgb]{0.02,0.16,0.49}{#1}}
\newcommand{\ImportTok}[1]{\textcolor[rgb]{0.00,0.50,0.00}{\textbf{#1}}}
\newcommand{\InformationTok}[1]{\textcolor[rgb]{0.38,0.63,0.69}{\textbf{\textit{#1}}}}
\newcommand{\KeywordTok}[1]{\textcolor[rgb]{0.00,0.44,0.13}{\textbf{#1}}}
\newcommand{\NormalTok}[1]{#1}
\newcommand{\OperatorTok}[1]{\textcolor[rgb]{0.40,0.40,0.40}{#1}}
\newcommand{\OtherTok}[1]{\textcolor[rgb]{0.00,0.44,0.13}{#1}}
\newcommand{\PreprocessorTok}[1]{\textcolor[rgb]{0.74,0.48,0.00}{#1}}
\newcommand{\RegionMarkerTok}[1]{#1}
\newcommand{\SpecialCharTok}[1]{\textcolor[rgb]{0.25,0.44,0.63}{#1}}
\newcommand{\SpecialStringTok}[1]{\textcolor[rgb]{0.73,0.40,0.53}{#1}}
\newcommand{\StringTok}[1]{\textcolor[rgb]{0.25,0.44,0.63}{#1}}
\newcommand{\VariableTok}[1]{\textcolor[rgb]{0.10,0.09,0.49}{#1}}
\newcommand{\VerbatimStringTok}[1]{\textcolor[rgb]{0.25,0.44,0.63}{#1}}
\newcommand{\WarningTok}[1]{\textcolor[rgb]{0.38,0.63,0.69}{\textbf{\textit{#1}}}}
\usepackage{longtable,booktabs,array}
\usepackage{calc} % for calculating minipage widths
% Correct order of tables after \paragraph or \subparagraph
\usepackage{etoolbox}
\makeatletter
\patchcmd\longtable{\par}{\if@noskipsec\mbox{}\fi\par}{}{}
\makeatother
% Allow footnotes in longtable head/foot
\IfFileExists{footnotehyper.sty}{\usepackage{footnotehyper}}{\usepackage{footnote}}
\makesavenoteenv{longtable}
\setlength{\emergencystretch}{3em} % prevent overfull lines
\providecommand{\tightlist}{%
  \setlength{\itemsep}{0pt}\setlength{\parskip}{0pt}}
\setcounter{secnumdepth}{-\maxdimen} % remove section numbering
% definitions for citeproc citations
\NewDocumentCommand\citeproctext{}{}
\NewDocumentCommand\citeproc{mm}{%
  \begingroup\def\citeproctext{#2}\cite{#1}\endgroup}
\makeatletter
 % allow citations to break across lines
 \let\@cite@ofmt\@firstofone
 % avoid brackets around text for \cite:
 \def\@biblabel#1{}
 \def\@cite#1#2{{#1\if@tempswa , #2\fi}}
\makeatother
\newlength{\cslhangindent}
\setlength{\cslhangindent}{1.5em}
\newlength{\csllabelwidth}
\setlength{\csllabelwidth}{3em}
\newenvironment{CSLReferences}[2] % #1 hanging-indent, #2 entry-spacing
 {\begin{list}{}{%
  \setlength{\itemindent}{0pt}
  \setlength{\leftmargin}{0pt}
  \setlength{\parsep}{0pt}
  % turn on hanging indent if param 1 is 1
  \ifodd #1
   \setlength{\leftmargin}{\cslhangindent}
   \setlength{\itemindent}{-1\cslhangindent}
  \fi
  % set entry spacing
  \setlength{\itemsep}{#2\baselineskip}}}
 {\end{list}}
\usepackage{calc}
\newcommand{\CSLBlock}[1]{\hfill\break\parbox[t]{\linewidth}{\strut\ignorespaces#1\strut}}
\newcommand{\CSLLeftMargin}[1]{\parbox[t]{\csllabelwidth}{\strut#1\strut}}
\newcommand{\CSLRightInline}[1]{\parbox[t]{\linewidth - \csllabelwidth}{\strut#1\strut}}
\newcommand{\CSLIndent}[1]{\hspace{\cslhangindent}#1}
\ifLuaTeX
  \usepackage{selnolig}  % disable illegal ligatures
\fi
\usepackage{bookmark}
\IfFileExists{xurl.sty}{\usepackage{xurl}}{} % add URL line breaks if available
\urlstyle{same}
\hypersetup{
  pdftitle={From File Labels to Byte Evidence: Budget-Aware Weakly Supervised Localization in Static PE Binaries},
  pdfauthor={{[}Author 1{]}; {[}Author 2{]}},
  hidelinks,
  pdfcreator={LaTeX via pandoc}}

\title{From File Labels to Byte Evidence: Budget-Aware Weakly Supervised
Localization in Static PE Binaries}
\usepackage{etoolbox}
\makeatletter
\providecommand{\subtitle}[1]{% add subtitle to \maketitle
  \apptocmd{\@title}{\par {\large #1 \par}}{}{}
}
\makeatother
\subtitle{UA-SAHI-Mal --- Venue-neutral pre-results manuscript draft}
\author{{[}Author 1{]} \and {[}Author 2{]}}
\date{Draft: 2026-09-04}

\begin{document}
\maketitle
\begin{abstract}
Static malware detectors can classify a Portable Executable (PE) file
without identifying the byte regions or functions that justify the
decision. This gap limits analyst trust and leaves reverse engineers
with most of the original program to inspect. Building supervised
localizers is difficult because public malware corpora predominantly
provide file-level labels, while independently validated component-level
annotations remain scarce, heterogeneous, or unavailable at scale. We
present UA-SAHI-Mal, a budget-aware weakly supervised framework that
learns from file labels to rank and recursively refine security-relevant
byte intervals in statically analyzable PE binaries. The framework
introduces PEAtlas, an invertible, structure-aware mapping between file
offsets, PE sections, virtual addresses, program-analysis entities, and
image coordinates. A relation-aware multiple-instance learner produces a
coarse evidence map; a structure-guided feature upsampler reconstructs
the map at higher resolution; and a guarded subset selector allocates a
fixed analysis budget to high-value and uncertain regions. The final
output is a ranked union of byte intervals, with masks and bounding
boxes derived only for visualization and detector compatibility. To
avoid evaluating localization solely against model explanations or a
single rule engine, we define a non-pooled, multi-tier benchmark
comprising controlled-provenance builds, rule-derived silver evidence,
source-grounded analyst-audited gold annotations, and an in-the-wild
analyst study. This pre-results manuscript specifies the threat model,
hypotheses, baselines, leakage controls, statistical analysis, and
reporting protocol. Experimental values are intentionally left as
{[}TBD{]} until measured.
\end{abstract}

{
\setcounter{tocdepth}{3}
\tableofcontents
}
\section{1. Introduction}\label{introduction}

Static machine-learning detectors have made it possible to score Windows
Portable Executable (PE) files at scale using engineered features, raw
bytes, or image-like representations
(\citeproc{ref-anderson2018ember}{Anderson and Roth 2018};
\citeproc{ref-raff2018malconv}{Raff et al. 2018};
\citeproc{ref-nataraj2011malwareimages}{Nataraj et al. 2011}). Their
standard output, however, is a file-level label or probability. A
malware analyst who receives a positive verdict must still identify the
functions, basic blocks, strings, embedded objects, or byte regions that
implement the behavior of interest. This distinction is operationally
important: a high-confidence classifier can reduce triage volume while
providing little guidance for reverse engineering.

Prior work has begun to address component localization. DeepReflect
ranks suspicious functions and basic blocks without requiring component
labels for training and showed that such prioritization can reduce
analyst workload (\citeproc{ref-downing2021deepreflect}{Downing et al.
2021}). Capability systems such as capa identify code-level behaviors
through expert-authored rules (\citeproc{ref-mandiant2026capa}{Mandiant
FLARE Team 2026}). More recent PE models use multiple-instance learning
(MIL) or position-preserving features to expose influential byte
segments (\citeproc{ref-peters2023highresolution}{Peters and Farhat
2023}; \citeproc{ref-kim2026pemilcon}{Kim, Trung, and Pham 2026};
\citeproc{ref-yevsikov2026dllicious}{Yevsikov and Nissim 2026}). These
systems establish that local evidence is useful, but they do not yield a
widely adopted public benchmark in which raw PE bytes, benign and
malicious files, independently validated malicious components,
reversible address mappings, and leakage-controlled splits are all
available together.

The absence of such a benchmark should not be interpreted as proof that
no component-level data exist, nor does it make supervised localization
impossible. Private corpora, manual reverse engineering, active
learning, and source-grounded builds are possible. The practical problem
is narrower: public, reproducible evaluation of malicious-component
localization remains constrained by annotation cost, licensing,
malware-handling restrictions, disassembly uncertainty, and ambiguity
about what constitutes a malicious component.

At the same time, weak supervision does not solve localization
automatically. A positive malware file contains substantial benign or
shared code, and a behavior may arise through interactions among
multiple functions. Attention weights or class activation maps can
reveal features used by a classifier, but they are not, by themselves,
evidence that a highlighted region implements malicious behavior.
High-resolution patch selection also has substantial precedent in MIL
and computational pathology (\citeproc{ref-ilse2018abmil}{Ilse, Tomczak,
and Welling 2018}; \citeproc{ref-lu2021clam}{Lu et al. 2021};
\citeproc{ref-li2021dsmil}{Li, Li, and Eliceiri 2021};
\citeproc{ref-shao2021transmil}{Shao et al. 2021}), and high-resolution
malware-image MIL has already used attention to prioritize patches
(\citeproc{ref-peters2023highresolution}{Peters and Farhat 2023}).
Consequently, a contribution based only on ``CAM or MIL followed by
top-\(K\) tiles'' would be insufficient.

This paper studies a more specific problem: \textbf{under an explicit
analysis budget, can a model trained primarily with file-level labels
retrieve byte intervals that independent evidence and analysts associate
with security-relevant malware components?} We introduce UA-SAHI-Mal, a
hierarchical evidence-retrieval framework for statically analyzable PE
files. The system first constructs an invertible, structure-aware
representation, PEAtlas. A relation-aware MIL model learns a coarse
tile-level evidence distribution. A structure-guided instance of
Upsample Anything (UA) (\citeproc{ref-seo2026upsample}{Seo, Hamilton,
and Kim 2026}) reconstructs higher-resolution evidence, and a guarded
budget selector balances high scores, uncertainty, and section coverage.
Selected regions are recursively refined until the budget is exhausted
or the desired byte granularity is reached.

The canonical output is a ranked set of byte intervals and mapped
program-analysis entities, not a two-dimensional bounding box. Bounding
boxes are derived views because a contiguous one-dimensional byte
interval can wrap across image rows, and optimized functions can occupy
multiple disjoint code chunks. This design keeps evaluation in the
native PE coordinate system while retaining compatibility with image
models and visual analyst interfaces.

We further introduce an evaluation protocol that separates evidence by
provenance rather than pooling incompatible labels.
Controlled-provenance builds test mapping and shortcut resistance. YARA
and capa matches provide scalable silver evidence. Source-grounded
builds and selected DeepReflect artifacts form an analyst-audited gold
set. Finally, a blinded analyst study evaluates whether ranked regions
reduce time and code reviewed. Each tier is reported separately.

This work makes the following intended contributions, subject to the
empirical gates defined in Section 8:

\begin{enumerate}
\def\labelenumi{\arabic{enumi}.}
\tightlist
\item
  \textbf{A reproducible audit and task definition.} We formalize public
  static-PE localization requirements, distinguish file detection from
  component retrieval, and document inclusion and exclusion criteria for
  existing datasets and methods.
\item
  \textbf{PEAtlas.} We provide an invertible mapping among raw file
  offsets, PE sections, RVAs, disassembly entities, and image
  coordinates, preserving disjoint byte intervals as the canonical
  annotation.
\item
  \textbf{A budget-aware weak localizer.} UA-SAHI-Mal combines
  relation-aware MIL, structure-guided evidence reconstruction,
  counterfactual and layout-consistency objectives, and guarded
  hierarchical selection under an explicit cost constraint.
\item
  \textbf{A non-pooled multi-tier benchmark.} We construct
  controlled-provenance, silver, source-grounded gold, and
  analyst-evaluation tiers with explicit provenance and leakage
  controls.
\item
  \textbf{Analyst-centered evaluation.} We report evidence recall at
  fixed byte, function, latency, and compute budgets, along with
  time-to-first-relevant-function and reviewed-code reduction.
\end{enumerate}

No numerical result is claimed in this draft. All values, sample counts,
and effect sizes marked \texttt{{[}TBD{]}} must be replaced only after
executing the registered protocol.

\section{2. Background, Motivation, and
Scope}\label{background-motivation-and-scope}

\subsection{2.1 Static PE malware
detection}\label{static-pe-malware-detection}

A PE file consists of headers, sections, optional overlays, and metadata
that define how on-disk bytes are mapped into memory. Large malware
datasets such as EMBER and EMBER2024 support file-level learning from
static features (\citeproc{ref-anderson2018ember}{Anderson and Roth
2018}; \citeproc{ref-joyce2025ember2024}{Joyce et al. 2025}). BODMAS
provides temporal labels and disarmed malware binaries, while its
original benign binaries are not publicly distributed
(\citeproc{ref-yang2021bodmas}{Yang et al. 2021}). These resources are
suitable for detection benchmarks but do not directly provide
malicious-component intervals.

Raw-byte neural networks avoid manual feature engineering, while
visualization methods reshape bytes into grayscale or multichannel
images (\citeproc{ref-raff2018malconv}{Raff et al. 2018};
\citeproc{ref-nataraj2011malwareimages}{Nataraj et al. 2011}). Image
representations enable the use of convolutional or vision-transformer
backbones, but they introduce layout choices. Adjacent pixels at a row
boundary may be far apart in program structure, and visual edges do not
necessarily correspond to semantic function boundaries. Our evaluation
therefore compares one-dimensional raw-byte, fixed-width raster, and
section-aligned representations.

\subsection{2.2 Local evidence and analyst
workload}\label{local-evidence-and-analyst-workload}

DeepReflect demonstrated that function prioritization can reduce
reverse-engineering effort and that unsupervised reconstruction can
identify malware components without component labels during training
(\citeproc{ref-downing2021deepreflect}{Downing et al. 2021}). capa
complements learned methods with expert-authored capability rules
operating at instruction, basic-block, function, and file scopes
(\citeproc{ref-mandiant2026capa}{Mandiant FLARE Team 2026}). YARA
provides exact offsets and lengths for matched strings or byte patterns
(\citeproc{ref-virustotal2026yara}{VirusTotal 2026}). These tools are
useful baselines and annotation sources, but their outputs have
different semantics:

\begin{itemize}
\tightlist
\item
  a YARA match is a signature interval;
\item
  a capa match is a capability-bearing program-analysis scope;
\item
  a model attribution is a region influential to a prediction;
\item
  a source-grounded annotation is a component linked to known
  implementation and analyst judgment.
\end{itemize}

We retain these distinctions throughout the paper.

\subsection{2.3 Weakly supervised
localization}\label{weakly-supervised-localization}

In MIL, a bag is labeled while individual instances are unlabelled
(\citeproc{ref-ilse2018abmil}{Ilse, Tomczak, and Welling 2018}).
Attention-based MIL can assign interpretable weights to instances, and
pathology systems such as CLAM, DSMIL, and TransMIL have localized
diagnostically relevant regions in gigapixel images using slide-level
labels (\citeproc{ref-lu2021clam}{Lu et al. 2021};
\citeproc{ref-li2021dsmil}{Li, Li, and Eliceiri 2021};
\citeproc{ref-shao2021transmil}{Shao et al. 2021}). Malware files create
a harder identification problem: positive bags may contain shared
libraries, runtime code, benign resources, and multiple cooperating
malicious functions. A file-level label does not imply that every tile
is positive, or even that one isolated tile is sufficient.

UA-SAHI-Mal treats benign bags as stronger negative supervision and
positive bags as positive-unlabeled mixtures. It uses cross-layout
consistency and counterfactual tests to reduce dependence on arbitrary
image placement. Even with these constraints, the output remains a
ranked hypothesis until validated against independent evidence.

\subsection{2.4 Slicing and feature
upsampling}\label{slicing-and-feature-upsampling}

SAHI applies an object detector to overlapping image slices and merges
detections, primarily to improve small-object detection
(\citeproc{ref-akyon2022sahi}{Akyon, Altinuc, and Temizel 2022}). Our
system is inspired by its resolution-preserving slicing principle but
differs in two ways. First, PE localization begins without dense object
labels. Second, the primary output is byte evidence rather than
natural-image objects. We therefore use the term \textbf{budgeted
hierarchical slicing} for the core retrieval mechanism. The name
``SAHI'' is retained in the system name to identify its design lineage;
any experiment described as ``SAHI'' uses an actual detector or
segmenter on slices and follows the original merge protocol.

Upsample Anything learns a per-input anisotropic Gaussian operator to
transfer low-resolution features or probability maps to high resolution
(\citeproc{ref-seo2026upsample}{Seo, Hamilton, and Kim 2026}). Its
edge-aware reconstruction is potentially useful, but raw-byte edges need
not be semantic. We evaluate byte-only and structure-guided variants and
charge all test-time optimization cost to the end-to-end budget.

\section{3. Threat Model and Task
Definition}\label{threat-model-and-task-definition}

\subsection{3.1 Scope}\label{scope}

The primary scope is statically analyzable native Windows PE32 and PE32+
files. The system does not execute samples. The main benchmark excludes,
or separately labels:

\begin{itemize}
\tightlist
\item
  packed samples that cannot be unpacked with the declared preprocessing
  pipeline;
\item
  self-modifying code;
\item
  managed .NET assemblies;
\item
  kernel drivers requiring specialized semantics;
\item
  files for which the selected disassembler cannot recover stable
  function boundaries.
\end{itemize}

Packed and optimized files are evaluated as domain-shift strata where
feasible, not silently mixed into the primary gold set.

\subsection{3.2 Adversary}\label{adversary}

The adversary may alter nonsemantic file layout, add padding or
overlays, reorder sections when functionality permits, use common
packers, or introduce benign-looking code. The initial study does not
claim robustness to an adaptive attacker with full knowledge of the
model and arbitrary semantics-preserving binary rewriting. We
nevertheless include padding, row-width, section-layout, compiler,
optimization, and selected packing perturbations to detect
representation shortcuts.

\subsection{3.3 Input and output}\label{input-and-output}

For a PE file \(F\) with on-disk byte sequence

\[
\mathbf{b} = (b_0, b_1, \ldots, b_{N-1}), \qquad b_o \in \{0,\ldots,255\},
\]

UA-SAHI-Mal receives the bytes and statically derived structural
metadata. It returns a ranked collection

\[
\mathcal{R}(F) = \{(I_k, s_k, m_k)\}_{k=1}^{K},
\]

where \(I_k\) is a union of one or more half-open file-offset intervals,
\(s_k\) is a confidence score, and \(m_k\) contains section, RVA, mapped
function/basic-block identifiers, and provenance metadata. The total
cost must satisfy

\[
C(\mathcal{R}) \le B,
\]

where \(B\) may be defined in bytes, tiles, functions, encoder calls,
FLOPs, or wall-clock time.

\subsection{3.4 Target semantics}\label{target-semantics}

The target is \textbf{security-relevant evidence associated with
malicious behavior}, not every byte in a malware file. For gold
annotations, analysts assign one of:

\begin{itemize}
\tightlist
\item
  \texttt{MALICIOUS\_COMPONENT}: directly implements or enables the
  malicious objective;
\item
  \texttt{SUPPORTING\_COMPONENT}: infrastructure required by the
  malicious component but not independently malicious;
\item
  \texttt{BENIGN\_OR\_SHARED}: runtime, library, or ordinary
  functionality;
\item
  \texttt{UNCERTAIN}: insufficient evidence for adjudication.
\end{itemize}

Primary gold localization metrics use \texttt{MALICIOUS\_COMPONENT};
sensitivity analyses may include supporting components. Silver evidence
is never relabeled as gold without analyst adjudication.

\section{4. Auditing the Ground-Truth
Gap}\label{auditing-the-ground-truth-gap}

\subsection{4.1 Audit question}\label{audit-question}

We ask whether a publicly accessible benchmark satisfies all of the
following:

\begin{enumerate}
\def\labelenumi{\arabic{enumi}.}
\tightlist
\item
  raw PE bytes or a provably reversible representation;
\item
  benign and malicious files;
\item
  independently validated malicious-component labels;
\item
  a mapping from labels to file offsets and program-analysis entities;
\item
  leakage-controlled family, temporal, source, and duplicate splits.
\end{enumerate}

The audit protocol, search strings, databases, dates, inclusion
criteria, and exclusion reasons will be released in Appendix A and the
artifact repository. We use the cautious conclusion ``no widely adopted
public benchmark was identified within the audit scope,'' rather than a
universal nonexistence claim.

\subsection{4.2 Position of existing
resources}\label{position-of-existing-resources}

\begin{longtable}[]{@{}
  >{\raggedright\arraybackslash}p{(\columnwidth - 10\tabcolsep) * \real{0.1500}}
  >{\raggedleft\arraybackslash}p{(\columnwidth - 10\tabcolsep) * \real{0.2000}}
  >{\raggedleft\arraybackslash}p{(\columnwidth - 10\tabcolsep) * \real{0.2000}}
  >{\raggedright\arraybackslash}p{(\columnwidth - 10\tabcolsep) * \real{0.1500}}
  >{\raggedright\arraybackslash}p{(\columnwidth - 10\tabcolsep) * \real{0.1500}}
  >{\raggedright\arraybackslash}p{(\columnwidth - 10\tabcolsep) * \real{0.1500}}@{}}
\toprule\noalign{}
\begin{minipage}[b]{\linewidth}\raggedright
Resource
\end{minipage} & \begin{minipage}[b]{\linewidth}\raggedleft
Raw/reversible PE
\end{minipage} & \begin{minipage}[b]{\linewidth}\raggedleft
Benign + malicious
\end{minipage} & \begin{minipage}[b]{\linewidth}\raggedright
Local evidence
\end{minipage} & \begin{minipage}[b]{\linewidth}\raggedright
Evidence status
\end{minipage} & \begin{minipage}[b]{\linewidth}\raggedright
Intended use
\end{minipage} \\
\midrule\noalign{}
\endhead
\bottomrule\noalign{}
\endlastfoot
EMBER / EMBER2024 & Features and metadata; raw files not generally
included & Yes & No component intervals & File labels & Detection and
temporal baselines \\
BODMAS & Disarmed malware binaries; benign originals not public & Labels
include both & No component intervals & File labels & Detection and time
split \\
DeepReflect artifacts & Limited malware/source artifacts & Limited &
Function/basic-block evidence & Source/analyst-grounded seed & Gold seed
and localization baseline \\
EMBER2024-capa & Function bytes and disassembly for many malicious
functions & No balanced raw-file corpus & Capability-bearing functions &
Silver & Pretraining and silver evaluation \\
YARA & Applied to local corpus & Depends on corpus & Exact match offsets
& Silver signature & Rule agreement \\
capa & Applied to local corpus & Depends on corpus &
Instruction/BB/function capabilities & Silver capability & Rule
agreement and baseline \\
Source-compiled corpus & Yes & Constructed controls & Source-linked
interval unions & Gold after audit & Primary controlled gold \\
\end{longtable}

The final paper will expand this table with citations and dataset
version identifiers.

\subsection{4.3 Why supervised detection remains a
baseline}\label{why-supervised-detection-remains-a-baseline}

The lack of a large public gold corpus does not make supervised learning
invalid. Once Tier 3 annotations exist, we train supervised
detector/segmenter baselines on the available gold training portion, use
cross-validation where necessary, and report uncertainty. This
comparison prevents weak supervision from receiving an artificial
advantage.

\section{5. PEAtlas: Invertible Structure-Preserving
Representation}\label{peatlas-invertible-structure-preserving-representation}

\subsection{5.1 Coordinate spaces}\label{coordinate-spaces}

PEAtlas distinguishes five coordinate systems:

\begin{enumerate}
\def\labelenumi{\arabic{enumi}.}
\tightlist
\item
  file offsets on disk;
\item
  RVAs in the loaded image;
\item
  VAs after adding the image base;
\item
  instruction, basic-block, and function entities recovered by analysis;
\item
  pixels or patches in a model representation.
\end{enumerate}

A PE section records, among other fields, \texttt{VirtualAddress},
\texttt{VirtualSize}, \texttt{PointerToRawData}, and
\texttt{SizeOfRawData}; raw and virtual extents may differ
(\citeproc{ref-microsoft2025pe}{Microsoft 2025b}). Therefore, every
conversion can fail explicitly:

\[
\operatorname{rva2off}(r) \in \{o, \bot_{\text{unmapped}}\},
\]

\[
\operatorname{off2rva}(o) \in \{r, \bot_{\text{nonloaded}}\}.
\]

Certificate data, debug data, overlays, and zero-filled virtual ranges
are tagged rather than coerced into a false one-to-one mapping.

\subsection{5.2 Fixed-width raster}\label{fixed-width-raster}

For a baseline image width \(W\), a valid file offset \(o\) maps to

\[
\pi_W(o) = (x,y) = (o \bmod W, \lfloor o/W \rfloor).
\]

The inverse is \(o=yW+x\) for valid pixels. Each pixel retains its
original offset. A byte interval crossing a row boundary becomes
multiple horizontal runs, not one enclosing rectangle.

\subsection{5.3 Section-aligned atlas}\label{section-aligned-atlas}

The proposed layout assigns each section and nonsection region to a
distinct band. Alignment pixels map to \(\bot\) and are excluded from
loss and metrics. The tensor includes:

\[
X = [v, e, q, r, w, x, m_{raw}, m_{loaded}],
\]

where \(v\) is normalized byte value, \(e\) local entropy, \(q\) a
section embedding or ID channel, \(r/w/x\) are section permissions, and
masks identify valid raw and memory-mapped bytes.

The atlas stores an explicit pixel-to-offset lookup table and an
interval-to-mask routine. We release round-trip tests for at least
\texttt{{[}TBD:\ 10\^{}6\ or\ more{]}} randomly sampled offsets and edge
cases.

\subsection{5.4 Program-analysis
mapping}\label{program-analysis-mapping}

A disassembler returns instruction, basic-block, and function address
sets. PEAtlas maps each recovered address to a set of file-offset
intervals when possible. A function may map to disjoint intervals
because of inlining, cold-code splitting, or compiler/linker
optimization. Optimized code can eliminate, merge, or inline functions,
and identical COMDAT folding can merge code
(\citeproc{ref-microsoft2021optimized}{Microsoft 2021},
\citeproc{ref-microsoft2025opt}{2025a}). Consequently, the annotation
schema stores interval unions rather than a single start/end pair.

\section{6. UA-SAHI-Mal}\label{ua-sahi-mal}

\subsection{6.1 Overview}\label{overview}

UA-SAHI-Mal comprises five stages:

\begin{enumerate}
\def\labelenumi{\arabic{enumi}.}
\tightlist
\item
  PEAtlas construction;
\item
  coarse tile encoding and relation-aware MIL;
\item
  structure-guided evidence reconstruction;
\item
  guarded subset selection under budget;
\item
  recursive fine localization and evidence export.
\end{enumerate}

A system diagram will appear as Figure 1.

\subsection{6.2 Tile encoder and relation-aware
MIL}\label{tile-encoder-and-relation-aware-mil}

PEAtlas is divided into possibly overlapping coarse tiles \(t_j\) with
positional and section metadata \(u_j\). An encoder produces

\[
h_j = f_\theta(t_j).
\]

A relation-aware aggregator, instantiated with ABMIL, DSMIL, or a
transformer during ablation, computes context-dependent evidence scores

\[
a_j = g_\phi(h_j, u_j, \{h_\ell,u_\ell\}_{\ell \ne j})
\]

and a file representation

\[
z = \sum_j \operatorname{softmax}(a)_j h_j.
\]

The file probability is \(p=\sigma(q(z))\).

For a benign bag, all tiles receive negative-instance regularization.
For a malicious bag, tile labels are unobserved and are treated as
positive-unlabeled. The model is not allowed to infer gold labels from
YARA or capa during the file-label-only training condition.

\subsection{6.3 Coarse evidence map}\label{coarse-evidence-map}

Tile scores are projected into a low-resolution evidence map \(E_{lr}\).
Overlap is combined using a calibrated aggregation rule. We compare
mean, maximum, noisy-OR, and learned aggregation. Scores are
temperature-calibrated on validation data without gold test annotations.

\subsection{6.4 Structure-guided UA}\label{structure-guided-ua}

We optimize a per-file upsampling operator using PEAtlas guidance \(G\):

\[
\phi^* = \arg\min_\phi \mathcal{L}_{guide}(U_\phi(G_{lr},G_{hr}),G_{hr}),
\]

then transfer the operator to evidence:

\[
E_{hr}=U_{\phi^*}(E_{lr},G_{hr}).
\]

We test two guidance conditions:

\begin{itemize}
\tightlist
\item
  byte-only grayscale guidance;
\item
  structural guidance using byte, entropy, section, permissions, and
  validity masks.
\end{itemize}

Nearest, bilinear, bicubic, guided/JBU, and UA are compared. UA
optimization time and memory are included in all online efficiency
numbers.

\subsection{6.5 Guarded budget selector}\label{guarded-budget-selector}

For candidate tile \(j\), the selector computes

\[
s_j = \hat e_j + \alpha \hat u_j + \beta \hat n_j + \gamma \hat c_j,
\]

where \(\hat e_j\) is evidence, \(\hat u_j\) predictive or ensemble
uncertainty, \(\hat n_j\) diversity/novelty relative to selected tiles,
and \(\hat c_j\) coverage value for underrepresented executable sections
or spatial regions.

We select \(S\) subject to

\[
\max_{S \subseteq \mathcal{T}} \sum_{j\in S}s_j, \qquad
\text{s.t. } \sum_{j\in S}c_j \le B.
\]

A differentiable soft-selection relaxation may be used during training;
inference uses a deterministic algorithm with fixed tie breaking. The
paper reports both tile-count and wall-clock budgets.

\subsection{6.6 Hierarchical refinement}\label{hierarchical-refinement}

Selected tiles are subdivided and re-encoded at higher resolution.
Refinement continues until one of the following holds:

\begin{itemize}
\tightlist
\item
  the minimum interval granularity is reached;
\item
  evidence concentration exceeds threshold \(\tau_e\);
\item
  uncertainty falls below \(\tau_u\);
\item
  the remaining budget cannot fund another split.
\end{itemize}

The fine stage produces a score for each retained subinterval.
Neighboring intervals may be merged only when the merge does not bridge
unrelated invalid or low-score bytes beyond a declared tolerance.

\subsection{6.7 Training objective}\label{training-objective}

The complete objective is

\[
\mathcal{L} = \mathcal{L}_{file}
+ \lambda_{neg}\mathcal{L}_{neg}
+ \lambda_{cons}\mathcal{L}_{layout}
+ \lambda_{cf}\mathcal{L}_{cf}
+ \lambda_{stab}\mathcal{L}_{stab}
+ \lambda_{cost}\mathcal{L}_{cost}.
\]

\subsubsection{File classification}\label{file-classification}

\[
\mathcal{L}_{file}=\operatorname{BCE}(p,y).
\]

\subsubsection{Negative-instance
regularization}\label{negative-instance-regularization}

For benign files,

\[
\mathcal{L}_{neg}=\frac{1}{|\mathcal{T}|}\sum_j -\log(1-\sigma(a_j)).
\]

\subsubsection{Cross-layout consistency}\label{cross-layout-consistency}

Let \(T_1\) and \(T_2\) be invertible layout transformations, such as
different raster widths or raster versus section atlas. After mapping
evidence back to offsets,

\[
\mathcal{L}_{layout}=D(E^{off}_{T_1},E^{off}_{T_2}).
\]

\subsubsection{Counterfactual sufficiency and
comprehensiveness}\label{counterfactual-sufficiency-and-comprehensiveness}

Let \(F_S\) retain selected bytes/tiles with a neutralized context, and
\(F_{\setminus S}\) neutralize selected evidence while preserving the
rest. We encourage

\[
p(F_S) \approx p(F), \qquad p(F_{\setminus S}) < p(F)-\delta.
\]

Controls verify that neutralization does not corrupt PE parsing or
introduce an obvious marker.

\subsubsection{Stability}\label{stability}

Functionality-preserving or representation-preserving perturbations
\(A(F)\) include permitted padding, row-width changes, and benign
overlay modifications:

\[
\mathcal{L}_{stab}=D(E^{off}(F),E^{off}(A(F))).
\]

\subsubsection{Cost}\label{cost}

\[
\mathcal{L}_{cost}=\max(0,C(S)-B).
\]

\section{7. Multi-Tier Evaluation
Benchmark}\label{multi-tier-evaluation-benchmark}

The tiers differ in semantic strength and are never pooled into a single
localization score.

\subsection{7.1 Tier 0: Mapping tests}\label{tier-0-mapping-tests}

Tier 0 verifies PEAtlas rather than machine learning. Tests include:

\begin{itemize}
\tightlist
\item
  file-offset/pixel round trips;
\item
  RVA/file-offset conversion for raw, zero-filled, and unmapped ranges;
\item
  section alignment and overlay handling;
\item
  interval-to-mask-to-interval exact recovery;
\item
  agreement among at least two parsers on supported fields;
\item
  disjoint function chunk representation.
\end{itemize}

A release gate requires zero unexplained mapping errors on
\texttt{{[}TBD{]}} deterministic and randomized cases.

\subsection{7.2 Tier 1: Controlled-provenance
builds}\label{tier-1-controlled-provenance-builds}

Tier 1 consists of benign host programs linked with safe, source-known
modules. Modules emulate security-relevant code patterns without harmful
external effects. Examples may include mock process enumeration, local
file transformation in a sandbox path, or network-client code redirected
to a local test service. Each module has a matched inert control with
similar size, entropy, compiler, and placement.

The build system randomizes:

\begin{itemize}
\tightlist
\item
  source module;
\item
  compiler and version;
\item
  architecture;
\item
  optimization mode;
\item
  section placement;
\item
  interval length;
\item
  padding and overlay conditions;
\item
  image layout width.
\end{itemize}

Ground truth is the union of source-linked code and data intervals
produced by debug and map information, verified against disassembly.
Tier 1 measures provenance localization and shortcut resistance, not
real-world malware accuracy. \texttt{secml\_malware} manipulations may
be used to generate padding/slack perturbations, but the library is not
treated as a malicious-payload annotator
(\citeproc{ref-demetrio2021secmlmalware}{Demetrio and Biggio 2021}).

\subsection{7.3 Tier 2: Rule-derived silver
evidence}\label{tier-2-rule-derived-silver-evidence}

Tier 2 applies YARA and capa to a legally authorized local PE corpus.
YARA provides match offsets and lengths; capa provides matched
instruction, basic-block, function, or file scopes
(\citeproc{ref-virustotal2026yara}{VirusTotal 2026};
\citeproc{ref-mandiant2026capa}{Mandiant FLARE Team 2026}).
EMBER2024-capa provides function-level bytes and disassembly for a large
collection of malicious-file functions and is used for fragment
pretraining or silver analysis, not whole-file coordinate reconstruction
unless the corresponding authorized PE is available
(\citeproc{ref-joyce2026embercapa}{Joyce and collaborators 2026}).

We report:

\begin{itemize}
\tightlist
\item
  overlap with YARA intervals;
\item
  capa function recall@budget;
\item
  capability-stratified results;
\item
  rule-family and capability holdout;
\item
  disagreement between learned evidence and rules.
\end{itemize}

No silver record is included in gold metrics without manual review.

\subsection{7.4 Tier 3: Source-grounded analyst-audited
gold}\label{tier-3-source-grounded-analyst-audited-gold}

Tier 3 combines selected DeepReflect source-grounded artifacts with
reproducible source builds. Each project is built in at least two modes:

\begin{enumerate}
\def\labelenumi{\arabic{enumi}.}
\tightlist
\item
  \textbf{Gold mode:} optimization and inlining disabled, no
  identical-code folding, symbols and map data retained;
\item
  \textbf{Shift mode:} realistic optimization enabled to evaluate
  transfer and annotation fragmentation.
\end{enumerate}

A source-to-binary pipeline yields candidate interval unions. Two
analysts independently assign component labels and supporting evidence.
Disagreements are adjudicated by a third reviewer or joint review. The
released schema records:

\begin{Shaded}
\begin{Highlighting}[]
\NormalTok{sample\_sha256}
\NormalTok{project\_id}
\NormalTok{build\_id}
\NormalTok{architecture}
\NormalTok{compiler\_and\_flags}
\NormalTok{packer\_status}
\NormalTok{component\_id}
\NormalTok{component\_semantics}
\NormalTok{label\_class}
\NormalTok{file\_offset\_intervals}
\NormalTok{rva\_intervals}
\NormalTok{mapped\_functions}
\NormalTok{source\_locations}
\NormalTok{annotator\_ids\_pseudonymous}
\NormalTok{agreement\_status}
\NormalTok{provenance}
\end{Highlighting}
\end{Shaded}

The gold set excludes ambiguous components from primary metrics and
reports them separately.

\subsection{7.5 Tier 4: In-the-wild analyst
audit}\label{tier-4-in-the-wild-analyst-audit}

A blinded crossover study compares UA-SAHI-Mal, DeepReflect, capa, and a
no-guidance control. Analysts receive the same time or function budget
and may use the same reverse-engineering environment. The presentation
order is randomized.

Primary outcomes:

\begin{itemize}
\tightlist
\item
  time to first relevant function;
\item
  number of functions reviewed before first relevant evidence;
\item
  precision among the first \(K\) recommendations;
\item
  fraction of analyst-confirmed relevant functions recovered;
\item
  analyst confidence and perceived usefulness.
\end{itemize}

The study protocol, recruitment criteria, consent, and compensation are
reviewed by the applicable institutional process before data collection.

\section{8. Experimental Methodology}\label{experimental-methodology}

\subsection{8.1 Pre-registration and empirical
gates}\label{pre-registration-and-empirical-gates}

Before evaluating the test sets, we freeze:

\begin{itemize}
\tightlist
\item
  primary metrics and budgets;
\item
  dataset split hashes;
\item
  gold and silver label versions;
\item
  model-selection metric;
\item
  hyperparameter search space;
\item
  exclusion rules;
\item
  analyst-study endpoints.
\end{itemize}

The project proceeds through five gates:

\begin{itemize}
\tightlist
\item
  \textbf{G0 Mapping integrity:} no unexplained coordinate errors;
\item
  \textbf{G1 Gold feasibility:} diverse, analyst-verifiable component
  intervals can be constructed;
\item
  \textbf{G2 Localization signal:} the model exceeds random, entropy,
  and base MIL rankings on validation gold;
\item
  \textbf{G3 Budget Pareto:} the method improves evidence recall at
  matched cost or reduces cost at matched recall;
\item
  \textbf{G4 Analyst utility:} the blinded study shows meaningful
  workload reduction;
\item
  \textbf{G5 Robustness:} evidence is not dominated by layout, padding,
  or compiler shortcuts.
\end{itemize}

Failure of G1 changes the claim from malicious-component localization to
weak evidence retrieval. Failure of G3 removes the efficiency claim.
Failure of G4 limits the paper to a benchmark/method study rather than
an analyst-productivity claim.

\subsection{8.2 Datasets}\label{datasets}

The final dataset table will report exact counts after legal and
technical validation.

\begin{longtable}[]{@{}
  >{\raggedright\arraybackslash}p{(\columnwidth - 12\tabcolsep) * \real{0.1200}}
  >{\raggedright\arraybackslash}p{(\columnwidth - 12\tabcolsep) * \real{0.1200}}
  >{\raggedleft\arraybackslash}p{(\columnwidth - 12\tabcolsep) * \real{0.1600}}
  >{\raggedleft\arraybackslash}p{(\columnwidth - 12\tabcolsep) * \real{0.1600}}
  >{\raggedleft\arraybackslash}p{(\columnwidth - 12\tabcolsep) * \real{0.1600}}
  >{\raggedleft\arraybackslash}p{(\columnwidth - 12\tabcolsep) * \real{0.1600}}
  >{\raggedright\arraybackslash}p{(\columnwidth - 12\tabcolsep) * \real{0.1200}}@{}}
\toprule\noalign{}
\begin{minipage}[b]{\linewidth}\raggedright
Split/source
\end{minipage} & \begin{minipage}[b]{\linewidth}\raggedright
Role
\end{minipage} & \begin{minipage}[b]{\linewidth}\raggedleft
Files
\end{minipage} & \begin{minipage}[b]{\linewidth}\raggedleft
Families/projects
\end{minipage} & \begin{minipage}[b]{\linewidth}\raggedleft
Gold components
\end{minipage} & \begin{minipage}[b]{\linewidth}\raggedleft
Silver components
\end{minipage} & \begin{minipage}[b]{\linewidth}\raggedright
Time range
\end{minipage} \\
\midrule\noalign{}
\endhead
\bottomrule\noalign{}
\endlastfoot
File-label training corpus & Weak training & {[}TBD{]} & {[}TBD{]} & 0 &
0 or excluded & {[}TBD{]} \\
Tier 1 controlled provenance & Sanity/robustness & {[}TBD{]} & {[}TBD{]}
& {[}TBD{]} & 0 & generated \\
Tier 2 rule evidence & Scalable silver & {[}TBD{]} & {[}TBD{]} & 0 &
{[}TBD{]} & {[}TBD{]} \\
Tier 3 source-grounded & Primary gold & {[}TBD{]} & {[}TBD{]} &
{[}TBD{]} & optional & {[}TBD{]} \\
Tier 4 analyst study & Utility & {[}TBD{]} & {[}TBD{]} &
analyst-confirmed & optional & {[}TBD{]} \\
\end{longtable}

EMBER/EMBER2024 features may support file-level baselines, but
whole-file image training requires authorized raw PEs. BODMAS disarmed
malware binaries may contribute malicious images, while benign raw
binaries require a separately licensed corpus
(\citeproc{ref-joyce2025ember2024}{Joyce et al. 2025};
\citeproc{ref-yang2021bodmas}{Yang et al. 2021}).

\subsection{8.3 Splitting and leakage
prevention}\label{splitting-and-leakage-prevention}

We construct grouped splits using:

\begin{itemize}
\tightlist
\item
  chronological collection time;
\item
  malware family;
\item
  source project;
\item
  compiler/toolchain;
\item
  packer;
\item
  exact file hash;
\item
  near-duplicate clustering;
\item
  function hash;
\item
  source-function identity across builds;
\item
  YARA/capa rule and capability group.
\end{itemize}

All builds of the same source component remain in one split unless the
experiment explicitly studies build transfer. EMBER2024-capa function
duplicates are removed or grouped before training and evaluation.

\subsection{8.4 Baselines}\label{baselines}

\subsubsection{File-level detection}\label{file-level-detection}

\begin{itemize}
\tightlist
\item
  EMBER LightGBM;
\item
  MalConv or MalConv2;
\item
  full-image CNN;
\item
  vision transformer;
\item
  Peters--Farhat high-resolution malware MIL;
\item
  PE-MILCon;
\item
  UA-SAHI-Mal file head.
\end{itemize}

\subsubsection{Localization and ranking}\label{localization-and-ranking}

\begin{itemize}
\tightlist
\item
  random-\(K\);
\item
  uniform grid/section-\(K\);
\item
  entropy-\(K\);
\item
  Grad-CAM and a higher-resolution CAM variant;
\item
  Integrated Gradients;
\item
  ABMIL;
\item
  CLAM;
\item
  DSMIL;
\item
  TransMIL;
\item
  Peters--Farhat attention;
\item
  PE-MILCon attention;
\item
  DeepReflect;
\item
  capa and YARA;
\item
  exhaustive tile scan;
\item
  supervised detector/segmenter trained on available Tier 3 gold;
\item
  DLLicious/X-MinHash-style attribution if a faithful implementation is
  available.
\end{itemize}

\subsubsection{Upsampling}\label{upsampling}

\begin{itemize}
\tightlist
\item
  nearest;
\item
  bilinear;
\item
  bicubic;
\item
  guided/JBU;
\item
  UA byte-only;
\item
  UA structural guidance.
\end{itemize}

\subsection{8.5 Metrics}\label{metrics}

\subsubsection{File-level metrics}\label{file-level-metrics}

\begin{itemize}
\tightlist
\item
  PR-AUC;
\item
  ROC-AUC;
\item
  F1;
\item
  TPR at FPR 0.1\% and 1\%;
\item
  expected calibration error.
\end{itemize}

\subsubsection{Localization metrics}\label{localization-metrics}

Let \(G\) be the union of gold byte intervals and \(R_B\) the retrieved
intervals under budget \(B\).

Byte recall:

\[
\operatorname{Recall}_{byte}(B)=\frac{|R_B\cap G|}{|G|}.
\]

Byte precision:

\[
\operatorname{Precision}_{byte}(B)=\frac{|R_B\cap G|}{|R_B|}.
\]

Function recall:

\[
\operatorname{Recall}_{func}(K)=\frac{\#\text{gold functions intersecting top-}K}{\#\text{gold functions}}.
\]

We additionally report:

\begin{itemize}
\tightlist
\item
  byte-level AUPRC;
\item
  interval IoU;
\item
  first relevant function rank;
\item
  reviewed-byte fraction;
\item
  reviewed-function fraction;
\item
  derived mask IoU and bbox AP50/AP75 as secondary metrics.
\end{itemize}

\subsubsection{Efficiency metrics}\label{efficiency-metrics}

\begin{itemize}
\tightlist
\item
  coarse encoder time;
\item
  UA optimization time;
\item
  selector time;
\item
  refinement time;
\item
  online disassembly/parsing time;
\item
  end-to-end wall-clock latency;
\item
  throughput;
\item
  peak RAM and VRAM;
\item
  encoder/detector calls;
\item
  FLOPs where meaningful.
\end{itemize}

Preprocessing is reported as online or amortized offline rather than
omitted.

\subsection{8.6 Statistical analysis}\label{statistical-analysis}

We use at least \texttt{{[}TBD{]}} independent seeds or bootstrap
resamples. Confidence intervals are computed at the file or project
group level to avoid treating correlated functions as independent.
Comparisons use paired bootstrap or permutation tests and report effect
sizes. Tier 1, Tier 2, Tier 3, and Tier 4 results are never averaged
into one score.

\subsection{8.7 Research questions}\label{research-questions}

\begin{itemize}
\tightlist
\item
  \textbf{RQ1 --- File-level validity:} Does weak evidence learning
  preserve competitive malware detection?
\item
  \textbf{RQ2 --- Budgeted localization:} At matched cost, does
  UA-SAHI-Mal retrieve more gold malicious components than attribution,
  MIL, rules, and DeepReflect?
\item
  \textbf{RQ3 --- Reconstruction:} Does structure-guided UA improve
  offset-level localization enough to justify its cost?
\item
  \textbf{RQ4 --- Objective and guards:} Which constraints reduce
  false-negative components and representation shortcuts?
\item
  \textbf{RQ5 --- Analyst utility:} Does ranked evidence reduce time and
  code reviewed?
\item
  \textbf{RQ6 --- Robustness:} How does localization change across
  layout, compiler, optimization, padding, family, temporal, and packing
  shifts?
\item
  \textbf{RQ7 --- Component scale:} Conditional on the observed
  gold-size distribution, are small components disproportionately missed
  by full-image models?
\end{itemize}

\section{9. Results}\label{results}

This section is a reporting template. It must not be populated with
estimated or illustrative values presented as measurements.

\subsection{9.1 RQ1: File-level
detection}\label{rq1-file-level-detection}

\textbf{Planned statement:} UA-SAHI-Mal achieved \texttt{{[}TBD{]}}
PR-AUC and \texttt{{[}TBD{]}} TPR at 0.1\% FPR on the temporal test set.
Relative to \texttt{{[}best\ baseline{]}}, the difference was
\texttt{{[}TBD{]}} with a 95\% confidence interval of
\texttt{{[}TBD{]}}.

\begin{longtable}[]{@{}lrrrrr@{}}
\toprule\noalign{}
Method & PR-AUC & ROC-AUC & TPR@0.1\% FPR & F1 & ECE \\
\midrule\noalign{}
\endhead
\bottomrule\noalign{}
\endlastfoot
EMBER LightGBM & TBD & TBD & TBD & TBD & TBD \\
MalConv & TBD & TBD & TBD & TBD & TBD \\
Full-image CNN & TBD & TBD & TBD & TBD & TBD \\
Peters MIL & TBD & TBD & TBD & TBD & TBD \\
PE-MILCon & TBD & TBD & TBD & TBD & TBD \\
UA-SAHI-Mal & TBD & TBD & TBD & TBD & TBD \\
\end{longtable}

\subsection{9.2 RQ2: Gold evidence
retrieval}\label{rq2-gold-evidence-retrieval}

\begin{longtable}[]{@{}
  >{\raggedright\arraybackslash}p{(\columnwidth - 10\tabcolsep) * \real{0.1304}}
  >{\raggedleft\arraybackslash}p{(\columnwidth - 10\tabcolsep) * \real{0.1739}}
  >{\raggedleft\arraybackslash}p{(\columnwidth - 10\tabcolsep) * \real{0.1739}}
  >{\raggedleft\arraybackslash}p{(\columnwidth - 10\tabcolsep) * \real{0.1739}}
  >{\raggedleft\arraybackslash}p{(\columnwidth - 10\tabcolsep) * \real{0.1739}}
  >{\raggedleft\arraybackslash}p{(\columnwidth - 10\tabcolsep) * \real{0.1739}}@{}}
\toprule\noalign{}
\begin{minipage}[b]{\linewidth}\raggedright
Method
\end{minipage} & \begin{minipage}[b]{\linewidth}\raggedleft
Byte AUPRC
\end{minipage} & \begin{minipage}[b]{\linewidth}\raggedleft
Function Recall@10
\end{minipage} & \begin{minipage}[b]{\linewidth}\raggedleft
Recall@25\% byte budget
\end{minipage} & \begin{minipage}[b]{\linewidth}\raggedleft
First-hit rank
\end{minipage} & \begin{minipage}[b]{\linewidth}\raggedleft
Reviewed-byte fraction
\end{minipage} \\
\midrule\noalign{}
\endhead
\bottomrule\noalign{}
\endlastfoot
Random & TBD & TBD & TBD & TBD & TBD \\
Entropy & TBD & TBD & TBD & TBD & TBD \\
Grad-CAM & TBD & TBD & TBD & TBD & TBD \\
ABMIL & TBD & TBD & TBD & TBD & TBD \\
DSMIL & TBD & TBD & TBD & TBD & TBD \\
TransMIL & TBD & TBD & TBD & TBD & TBD \\
DeepReflect & TBD & TBD & TBD & TBD & TBD \\
Full scan & TBD & TBD & 1.000 & TBD & 1.000 \\
UA-SAHI-Mal & TBD & TBD & TBD & TBD & TBD \\
\end{longtable}

Figure 2 will plot gold component recall against end-to-end latency and
reviewed-byte fraction.

\subsection{9.3 RQ3: Upsampling}\label{rq3-upsampling}

\begin{longtable}[]{@{}
  >{\raggedright\arraybackslash}p{(\columnwidth - 10\tabcolsep) * \real{0.1304}}
  >{\raggedleft\arraybackslash}p{(\columnwidth - 10\tabcolsep) * \real{0.1739}}
  >{\raggedleft\arraybackslash}p{(\columnwidth - 10\tabcolsep) * \real{0.1739}}
  >{\raggedleft\arraybackslash}p{(\columnwidth - 10\tabcolsep) * \real{0.1739}}
  >{\raggedleft\arraybackslash}p{(\columnwidth - 10\tabcolsep) * \real{0.1739}}
  >{\raggedleft\arraybackslash}p{(\columnwidth - 10\tabcolsep) * \real{0.1739}}@{}}
\toprule\noalign{}
\begin{minipage}[b]{\linewidth}\raggedright
Reconstruction
\end{minipage} & \begin{minipage}[b]{\linewidth}\raggedleft
Byte AUPRC
\end{minipage} & \begin{minipage}[b]{\linewidth}\raggedleft
Interval IoU
\end{minipage} & \begin{minipage}[b]{\linewidth}\raggedleft
Layout stability
\end{minipage} & \begin{minipage}[b]{\linewidth}\raggedleft
Added latency
\end{minipage} & \begin{minipage}[b]{\linewidth}\raggedleft
Peak VRAM
\end{minipage} \\
\midrule\noalign{}
\endhead
\bottomrule\noalign{}
\endlastfoot
Nearest & TBD & TBD & TBD & TBD & TBD \\
Bilinear & TBD & TBD & TBD & TBD & TBD \\
Bicubic & TBD & TBD & TBD & TBD & TBD \\
JBU & TBD & TBD & TBD & TBD & TBD \\
UA, byte-only & TBD & TBD & TBD & TBD & TBD \\
UA, structural & TBD & TBD & TBD & TBD & TBD \\
\end{longtable}

The text will explicitly report whether UA remains Pareto-efficient
after charging its test-time optimization cost.

\subsection{9.4 RQ4: Ablation}\label{rq4-ablation}

\begin{longtable}[]{@{}
  >{\raggedright\arraybackslash}p{(\columnwidth - 10\tabcolsep) * \real{0.1304}}
  >{\raggedleft\arraybackslash}p{(\columnwidth - 10\tabcolsep) * \real{0.1739}}
  >{\raggedleft\arraybackslash}p{(\columnwidth - 10\tabcolsep) * \real{0.1739}}
  >{\raggedleft\arraybackslash}p{(\columnwidth - 10\tabcolsep) * \real{0.1739}}
  >{\raggedleft\arraybackslash}p{(\columnwidth - 10\tabcolsep) * \real{0.1739}}
  >{\raggedleft\arraybackslash}p{(\columnwidth - 10\tabcolsep) * \real{0.1739}}@{}}
\toprule\noalign{}
\begin{minipage}[b]{\linewidth}\raggedright
Variant
\end{minipage} & \begin{minipage}[b]{\linewidth}\raggedleft
File PR-AUC
\end{minipage} & \begin{minipage}[b]{\linewidth}\raggedleft
Gold function recall
\end{minipage} & \begin{minipage}[b]{\linewidth}\raggedleft
Byte AUPRC
\end{minipage} & \begin{minipage}[b]{\linewidth}\raggedleft
Latency
\end{minipage} & \begin{minipage}[b]{\linewidth}\raggedleft
Padding stability
\end{minipage} \\
\midrule\noalign{}
\endhead
\bottomrule\noalign{}
\endlastfoot
Full model & TBD & TBD & TBD & TBD & TBD \\
-- UA & TBD & TBD & TBD & TBD & TBD \\
-- counterfactual loss & TBD & TBD & TBD & TBD & TBD \\
-- layout consistency & TBD & TBD & TBD & TBD & TBD \\
-- uncertainty guard & TBD & TBD & TBD & TBD & TBD \\
-- coverage guard & TBD & TBD & TBD & TBD & TBD \\
byte-only representation & TBD & TBD & TBD & TBD & TBD \\
raster instead of section atlas & TBD & TBD & TBD & TBD & TBD \\
\end{longtable}

\subsection{9.5 RQ5: Analyst study}\label{rq5-analyst-study}

\begin{longtable}[]{@{}
  >{\raggedright\arraybackslash}p{(\columnwidth - 10\tabcolsep) * \real{0.1304}}
  >{\raggedleft\arraybackslash}p{(\columnwidth - 10\tabcolsep) * \real{0.1739}}
  >{\raggedleft\arraybackslash}p{(\columnwidth - 10\tabcolsep) * \real{0.1739}}
  >{\raggedleft\arraybackslash}p{(\columnwidth - 10\tabcolsep) * \real{0.1739}}
  >{\raggedleft\arraybackslash}p{(\columnwidth - 10\tabcolsep) * \real{0.1739}}
  >{\raggedleft\arraybackslash}p{(\columnwidth - 10\tabcolsep) * \real{0.1739}}@{}}
\toprule\noalign{}
\begin{minipage}[b]{\linewidth}\raggedright
Condition
\end{minipage} & \begin{minipage}[b]{\linewidth}\raggedleft
Median time-to-first evidence
\end{minipage} & \begin{minipage}[b]{\linewidth}\raggedleft
Functions reviewed
\end{minipage} & \begin{minipage}[b]{\linewidth}\raggedleft
Precision@10
\end{minipage} & \begin{minipage}[b]{\linewidth}\raggedleft
Confidence
\end{minipage} & \begin{minipage}[b]{\linewidth}\raggedleft
Workload score
\end{minipage} \\
\midrule\noalign{}
\endhead
\bottomrule\noalign{}
\endlastfoot
No guidance & TBD & TBD & TBD & TBD & TBD \\
capa & TBD & TBD & TBD & TBD & TBD \\
DeepReflect & TBD & TBD & TBD & TBD & TBD \\
UA-SAHI-Mal & TBD & TBD & TBD & TBD & TBD \\
\end{longtable}

The analysis will use participant-aware paired statistics and disclose
sample size, experience, order effects, and exclusions.

\subsection{9.6 RQ6: Robustness}\label{rq6-robustness}

We report stratified localization under:

\begin{itemize}
\tightlist
\item
  raster width changes;
\item
  benign padding and overlay insertion;
\item
  compiler and optimization changes;
\item
  family and time shifts;
\item
  packer strata where static analysis succeeds;
\item
  disassembler disagreement.
\end{itemize}

A model that preserves file classification but moves evidence to
inserted padding fails the localization robustness criterion.

\subsection{9.7 RQ7: Component scale}\label{rq7-component-scale}

We first report the empirical distribution of gold interval length,
function size, image-mask area, fragmentation, and post-resize pixel
count. Only if small components are sufficiently represented do we
report scale-stratified recall. COCO small-object thresholds are not
adopted without justification because PE byte intervals have different
semantics from natural-image objects.

\section{10. Security Analysis}\label{security-analysis}

\subsection{10.1 Shortcut tests}\label{shortcut-tests}

We test whether the model relies on:

\begin{itemize}
\tightlist
\item
  file size;
\item
  section count or names;
\item
  packer markers;
\item
  padding entropy;
\item
  certificate presence;
\item
  import-table artifacts;
\item
  family-specific build signatures;
\item
  image row boundaries.
\end{itemize}

Controls match these attributes while changing component provenance.
Cross-layout consistency is evaluated in offset space.

\subsection{10.2 Adaptive considerations}\label{adaptive-considerations}

A knowledgeable attacker may move or split malicious functionality,
interleave it with benign code, or add decoy high-scoring regions. We
evaluate bounded transformations and discuss the remaining adaptive
surface. The system is an analyst-prioritization aid, not a standalone
guarantee that unselected code is benign.

\subsection{10.3 Parser and disassembler
disagreement}\label{parser-and-disassembler-disagreement}

We compare at least two PE parsers and, on a subset, at least two
disassembly backends. Disagreement is recorded as annotation
uncertainty. Samples with unresolved mapping differences are excluded
from primary gold metrics under a preregistered rule and retained for
failure analysis.

\section{11. Case Studies}\label{case-studies}

Each case study will include:

\begin{enumerate}
\def\labelenumi{\arabic{enumi}.}
\tightlist
\item
  a PEAtlas view with ranked intervals;
\item
  original file offsets and mapped RVAs;
\item
  functions/basic blocks in the reverse-engineering tool;
\item
  source or analyst evidence supporting the label;
\item
  comparison with capa, DeepReflect, and one MIL baseline;
\item
  counterfactual deletion and layout-stability checks;
\item
  analysis of false positives and missed components.
\end{enumerate}

At least one success, one partial success, and one failure case will be
shown. Case studies will not replace aggregate evaluation.

\section{12. Discussion}\label{discussion}

\subsection{12.1 What localization means}\label{what-localization-means}

UA-SAHI-Mal retrieves regions associated with a file-level decision and
validated against available evidence. It does not prove causality in the
program-execution sense. A component can be security-relevant because it
enables persistence, evasion, command and control, credential access, or
another malicious objective, but static evidence may be incomplete or
ambiguous.

\subsection{12.2 Why not train a detector
directly?}\label{why-not-train-a-detector-directly}

We do train supervised baselines on available gold data. The weak
formulation is motivated by the scale mismatch between file labels and
component annotations, not by a claim that supervised detection is
conceptually invalid. The comparison quantifies when weak supervision is
useful and where it falls short.

\subsection{12.3 Why use an image
representation?}\label{why-use-an-image-representation}

Images provide a convenient multiscale spatial substrate and access to
mature visual encoders and upsamplers. However, the one-dimensional byte
domain remains authoritative. The representation comparison and
cross-layout objective test whether any benefit survives arbitrary
rasterization choices.

\subsection{12.4 Relationship to SAHI}\label{relationship-to-sahi}

The proposed selection mechanism is SAHI-inspired. We reserve the term
``Full SAHI'' for experiments that run an actual detector over all
slices and merge detections according to the SAHI protocol. Other
experiments are described as hierarchical slicing to avoid conflating
patch MIL with object detection.

\section{13. Ethics, Safety, and Artifact
Handling}\label{ethics-safety-and-artifact-handling}

The project handles malware only in isolated, access-controlled
environments. Samples are never executed on production systems. Raw
malware redistribution follows source licenses, provider agreements,
institutional policy, and applicable law. Where binaries cannot be
redistributed, the artifact releases hashes, build recipes, derived
annotations, feature records, and evaluation code to the extent
permitted.

Controlled-provenance modules are designed to avoid harmful external
effects. Analyst-study data are pseudonymized, and participants receive
informed consent and compensation under applicable institutional review.
The system is framed as decision support; analysts are warned that
unselected regions may still contain malicious behavior.

\section{14. Limitations}\label{limitations}

\begin{enumerate}
\def\labelenumi{\arabic{enumi}.}
\tightlist
\item
  Static analysis can miss runtime-decrypted, dynamically resolved,
  self-modifying, or environment-dependent behavior.
\item
  Source-grounded builds may differ from in-the-wild malware in
  optimization, packing, compiler, and operational context.
\item
  YARA and capa provide partial, rule-dependent evidence and cannot
  serve as complete maliciousness ground truth.
\item
  Debug information can describe code layout without deciding whether a
  function is malicious; analyst adjudication remains necessary.
\item
  File-level labels may not identify distributed multi-function
  behavior, even with MIL and counterfactual constraints.
\item
  UA may sharpen structural boundaries that are unrelated to semantic
  behavior and may not justify its online cost.
\item
  The primary gold corpus may remain small relative to file-level
  detection datasets, limiting family coverage.
\item
  An adaptive attacker can target the selector with decoys or
  distributed payloads.
\item
  Results for unpacked native PE files may not generalize to packed
  files, .NET, drivers, scripts, or other formats.
\end{enumerate}

\section{15. Related Work}\label{related-work}

\subsection{Static PE detection and malware
images}\label{static-pe-detection-and-malware-images}

Malware-image classification maps byte patterns to visual texture
(\citeproc{ref-nataraj2011malwareimages}{Nataraj et al. 2011}). Raw-byte
networks such as MalConv avoid image transformation
(\citeproc{ref-raff2018malconv}{Raff et al. 2018}), while EMBER provides
a strong engineered-feature baseline
(\citeproc{ref-anderson2018ember}{Anderson and Roth 2018}). EMBER2024
expands formats, tasks, and temporal evaluation
(\citeproc{ref-joyce2025ember2024}{Joyce et al. 2025}). UA-SAHI-Mal
differs by targeting independently evaluated component evidence under a
budget.

\subsection{Malware component localization and
explainability}\label{malware-component-localization-and-explainability}

DeepReflect localizes anomalous functions and incrementally clusters
malware functionality (\citeproc{ref-downing2021deepreflect}{Downing et
al. 2021}). capa matches expert-defined capabilities
(\citeproc{ref-mandiant2026capa}{Mandiant FLARE Team 2026}), and
DLLicious uses position-preserving attribution to indicate byte regions
associated with malicious DLL patterns
(\citeproc{ref-yevsikov2026dllicious}{Yevsikov and Nissim 2026}). These
are direct baselines or complementary evidence sources.

\subsection{Weak supervision and MIL}\label{weak-supervision-and-mil}

ABMIL introduced attention-based permutation-invariant aggregation
(\citeproc{ref-ilse2018abmil}{Ilse, Tomczak, and Welling 2018}). CLAM,
DSMIL, and TransMIL model representative and correlated regions in very
high-resolution weakly labeled images (\citeproc{ref-lu2021clam}{Lu et
al. 2021}; \citeproc{ref-li2021dsmil}{Li, Li, and Eliceiri 2021};
\citeproc{ref-shao2021transmil}{Shao et al. 2021}). Peters and Farhat
applied patch MIL to high-resolution malware images
(\citeproc{ref-peters2023highresolution}{Peters and Farhat 2023}), and
PE-MILCon combines byte-segment MIL with PE structural features
(\citeproc{ref-kim2026pemilcon}{Kim, Trung, and Pham 2026}). Our focus
is not attention visualization alone, but budgeted retrieval evaluated
against external evidence.

\subsection{High-resolution selective
inference}\label{high-resolution-selective-inference}

SAHI performs detector inference over slices to recover small objects
(\citeproc{ref-akyon2022sahi}{Akyon, Altinuc, and Temizel 2022}).
Upsample Anything reconstructs low-resolution features and probability
maps using per-input, edge-aware optimization
(\citeproc{ref-seo2026upsample}{Seo, Hamilton, and Kim 2026}).
UA-SAHI-Mal adapts these concepts to a weakly supervised byte-evidence
task and evaluates whether their added cost and representation
assumptions are justified.

\section{16. Conclusion}\label{conclusion}

This paper defines static PE malware localization as budgeted retrieval
of analyst-verifiable byte evidence rather than direct prediction of
unverified image boxes. UA-SAHI-Mal combines an invertible PE
representation, weak tile learning, structure-guided reconstruction,
guarded subset selection, and recursive refinement. Its evaluation
separates controlled provenance, silver rules, source-grounded gold, and
analyst utility. The design deliberately avoids treating model
attention, injected markers, or capability rules as interchangeable
ground truth. The empirical study will determine whether the proposed
system improves the evidence-recall--cost frontier and reduces
reverse-engineering workload. Until those experiments are complete, all
performance claims remain open hypotheses.

\section{Appendix A. Planned Systematic Audit
Protocol}\label{appendix-a.-planned-systematic-audit-protocol}

The final appendix will include:

\begin{itemize}
\tightlist
\item
  databases and repositories searched;
\item
  exact search strings;
\item
  search and update dates;
\item
  language and publication filters;
\item
  inclusion criteria;
\item
  duplicate handling;
\item
  two-reviewer screening procedure;
\item
  evidence extraction form;
\item
  reasons each dataset fails or satisfies the five benchmark conditions;
\item
  a public machine-readable audit table.
\end{itemize}

\section{Appendix B. Annotation
Schema}\label{appendix-b.-annotation-schema}

\begin{Shaded}
\begin{Highlighting}[]
\FunctionTok{\{}
  \DataTypeTok{"sample\_sha256"}\FunctionTok{:} \StringTok{"..."}\FunctionTok{,}
  \DataTypeTok{"build\_id"}\FunctionTok{:} \StringTok{"..."}\FunctionTok{,}
  \DataTypeTok{"coordinate\_version"}\FunctionTok{:} \StringTok{"peatlas{-}v1"}\FunctionTok{,}
  \DataTypeTok{"component\_id"}\FunctionTok{:} \StringTok{"..."}\FunctionTok{,}
  \DataTypeTok{"label\_class"}\FunctionTok{:} \StringTok{"MALICIOUS\_COMPONENT"}\FunctionTok{,}
  \DataTypeTok{"semantic\_tags"}\FunctionTok{:} \OtherTok{[}\StringTok{"..."}\OtherTok{]}\FunctionTok{,}
  \DataTypeTok{"file\_offset\_intervals"}\FunctionTok{:} \OtherTok{[[}\DecValTok{1000}\OtherTok{,} \DecValTok{1040}\OtherTok{],} \OtherTok{[}\DecValTok{2048}\OtherTok{,} \DecValTok{2112}\OtherTok{]]}\FunctionTok{,}
  \DataTypeTok{"rva\_intervals"}\FunctionTok{:} \OtherTok{[[}\DecValTok{8192}\OtherTok{,} \DecValTok{8232}\OtherTok{],} \OtherTok{[}\DecValTok{12288}\OtherTok{,} \DecValTok{12352}\OtherTok{]]}\FunctionTok{,}
  \DataTypeTok{"mapped\_functions"}\FunctionTok{:} \OtherTok{[}\StringTok{"sub\_401000"}\OtherTok{]}\FunctionTok{,}
  \DataTypeTok{"source\_locations"}\FunctionTok{:} \OtherTok{[}\StringTok{"module.c:20{-}71"}\OtherTok{]}\FunctionTok{,}
  \DataTypeTok{"provenance"}\FunctionTok{:} \StringTok{"SOURCE\_AND\_ANALYST"}\FunctionTok{,}
  \DataTypeTok{"annotator\_agreement"}\FunctionTok{:} \StringTok{"AGREED"}\FunctionTok{,}
  \DataTypeTok{"notes"}\FunctionTok{:} \StringTok{"..."}
\FunctionTok{\}}
\end{Highlighting}
\end{Shaded}

\section{Appendix C. Minimum Reproducibility
Checklist}\label{appendix-c.-minimum-reproducibility-checklist}

\begin{itemize}
\tightlist
\item
  source code and pinned environments;
\item
  dataset and split hashes;
\item
  PEAtlas mapping tests;
\item
  compiler/build containers for controlled and gold builds;
\item
  model checkpoints and hyperparameters;
\item
  all baseline adapters;
\item
  inference-time accounting script;
\item
  annotation guidelines and adjudication log;
\item
  statistical-analysis notebook;
\item
  failure-case inventory;
\item
  malware-handling and access instructions.
\end{itemize}

\phantomsection\label{refs}
\begin{CSLReferences}{1}{0}
\bibitem[\citeproctext]{ref-akyon2022sahi}
Akyon, Fatih Cagatay, Sinan Onur Altinuc, and Alptekin Temizel. 2022.
{``Slicing Aided Hyper Inference and Fine-Tuning for Small Object
Detection.''} In \emph{2022 IEEE International Conference on Image
Processing (ICIP)}, 966--70.
\url{https://doi.org/10.1109/ICIP46576.2022.9897990}.

\bibitem[\citeproctext]{ref-anderson2018ember}
Anderson, Hyrum S., and Phil Roth. 2018. {``{EMBER}: An Open Dataset for
Training Static PE Malware Machine Learning Models.''} \emph{arXiv
Preprint arXiv:1804.04637}. \url{https://arxiv.org/abs/1804.04637}.

\bibitem[\citeproctext]{ref-demetrio2021secmlmalware}
Demetrio, Luca, and Battista Biggio. 2021. {``Secml-Malware: A Python
Library for Adversarial Robustness Evaluation of Windows Malware
Classifiers.''} \emph{arXiv Preprint arXiv:2104.12848}.
\url{https://github.com/pralab/secml_malware}.

\bibitem[\citeproctext]{ref-downing2021deepreflect}
Downing, Evan, Yisroel Mirsky, Kyuhong Park, and Wenke Lee. 2021.
{``{DeepReflect}: Discovering Malicious Functionality Through Binary
Reconstruction.''} In \emph{30th USENIX Security Symposium (USENIX
Security 21)}, 3469--86. USENIX Association.
\url{https://www.usenix.org/conference/usenixsecurity21/presentation/downing}.

\bibitem[\citeproctext]{ref-ilse2018abmil}
Ilse, Maximilian, Jakub Tomczak, and Max Welling. 2018.
{``Attention-Based Deep Multiple Instance Learning.''} In
\emph{Proceedings of the 35th International Conference on Machine
Learning}, 80:2127--36. Proceedings of Machine Learning Research. PMLR.
\url{https://proceedings.mlr.press/v80/ilse18a.html}.

\bibitem[\citeproctext]{ref-joyce2026embercapa}
Joyce, Robert J. 2025. {``{EMBER2024-capa}:
Function-Level Capability, Byte, and Disassembly Records.''} Hugging
Face dataset.
\url{https://huggingface.co/datasets/joyce8/EMBER2024-capa}.

\bibitem[\citeproctext]{ref-joyce2025ember2024}
Joyce, Robert J., Gideon Miller, Phil Roth, Richard Zak, Elliott
Zaresky-Williams, Hyrum Anderson, Edward Raff, and James Holt. 2025.
{``{EMBER2024}: A Benchmark Dataset for Holistic Evaluation of Malware
Classifiers.''} In \emph{Proceedings of the 31st ACM SIGKDD Conference
on Knowledge Discovery and Data Mining}.
\url{https://doi.org/10.1145/3711896.3737431}.

\bibitem[\citeproctext]{ref-kim2026pemilcon}
Kim, Tuan Nguyen, Son Doan Trung, and Nguyen Minh Nhut Pham. 2026.
{``{PE-MILCon}: Multiple-Instance Learning with Contrastive Multi-View
Representation for Static Windows PE Malware Detection.''}
\emph{Computers, Materials \& Continua} 89 (1): 43.
\url{https://doi.org/10.32604/cmc.2026.084268}.

\bibitem[\citeproctext]{ref-li2021dsmil}
Li, Bin, Yin Li, and Kevin W. Eliceiri. 2021. {``Dual-Stream Multiple
Instance Learning Network for Whole Slide Image Classification with
Self-Supervised Contrastive Learning.''} In \emph{Proceedings of the
IEEE/CVF Conference on Computer Vision and Pattern Recognition}.
\url{https://arxiv.org/abs/2011.08939}.

\bibitem[\citeproctext]{ref-lu2021clam}
Lu, Ming Y., Drew F. K. Williamson, Tiffany Y. Chen, Richard J. Chen,
Matteo Barbieri, and Faisal Mahmood. 2021. {``Data-Efficient and Weakly
Supervised Computational Pathology on Whole-Slide Images.''}
\emph{Nature Biomedical Engineering} 5: 555--70.
\url{https://doi.org/10.1038/s41551-020-00682-w}.

\bibitem[\citeproctext]{ref-mandiant2026capa}
Mandiant FLARE Team. 2026. {``Capa: Extract Capabilities from Executable
Files.''} Software and documentation.
\url{https://mandiant.github.io/capa/}.

\bibitem[\citeproctext]{ref-microsoft2021optimized}
Microsoft. 2021. {``Debugging Optimized Code and Inline Functions.''}
\url{https://learn.microsoft.com/en-us/windows-hardware/drivers/debugger/debugging-optimized-code-and-inline-functions-external}.

\bibitem[\citeproctext]{ref-microsoft2025opt}
---------. 2025a. {``/OPT (Optimizations).''}
\url{https://learn.microsoft.com/en-us/cpp/build/reference/opt-optimizations}.

\bibitem[\citeproctext]{ref-microsoft2025pe}
---------. 2025b. {``PE Format.''}
\url{https://learn.microsoft.com/en-us/windows/win32/debug/pe-format}.

\bibitem[\citeproctext]{ref-nataraj2011malwareimages}
Nataraj, Lakshmanan, Sreejith Karthikeyan, Gregoire Jacob, and B. S.
Manjunath. 2011. {``Malware Images: Visualization and Automatic
Classification.''} In \emph{Proceedings of the 8th International
Symposium on Visualization for Cyber Security}.
\url{https://doi.org/10.1145/2016904.2016908}.

\bibitem[\citeproctext]{ref-peters2023highresolution}
Peters, Tim, and Hikmat Farhat. 2023. {``High-Resolution Image-Based
Malware Classification Using Multiple Instance Learning.''} \emph{arXiv
Preprint arXiv:2311.12760}.
\url{https://doi.org/10.48550/arXiv.2311.12760}.

\bibitem[\citeproctext]{ref-raff2018malconv}
Raff, Edward, Jon Barker, Jared Sylvester, Robert Brandon, Bryan
Catanzaro, and Charles Nicholas. 2018. {``Malware Detection by Eating a
Whole EXE.''} \emph{arXiv Preprint arXiv:1710.09435}.
\url{https://arxiv.org/abs/1710.09435}.

\bibitem[\citeproctext]{ref-seo2026upsample}
Seo, Minseok, Mark Hamilton, and Changick Kim. 2026. {``Upsample
Anything: A Simple and Hard to Beat Baseline for Feature Upsampling.''}
In \emph{Proceedings of the IEEE/CVF Conference on Computer Vision and
Pattern Recognition}. \url{https://arxiv.org/abs/2511.16301}.

\bibitem[\citeproctext]{ref-shao2021transmil}
Shao, Zhuchen, Hao Bian, Yang Chen, Yifeng Wang, Jian Zhang, Xiangyang
Ji, and Yongbing Zhang. 2021. {``{TransMIL}: Transformer-Based
Correlated Multiple Instance Learning for Whole Slide Image
Classification.''} In \emph{Advances in Neural Information Processing
Systems}. Vol. 34.
\url{https://proceedings.neurips.cc/paper/2021/hash/10c272d06794d3e5785d5e7c5356e9ff-Abstract.html}.

\bibitem[\citeproctext]{ref-virustotal2026yara}
VirusTotal. 2026. {``YARA c API Documentation.''}
\url{https://yara.readthedocs.io/en/latest/capi.html}.

\bibitem[\citeproctext]{ref-yang2021bodmas}
Yang, Limin, Arridhana Ciptadi, Ihar Laziuk, Ali Ahmadzadeh, and Gang
Wang. 2021. {``{BODMAS}: An Open Dataset for Learning-Based Temporal
Analysis of PE Malware.''} In \emph{4th Deep Learning and Security
Workshop}. \url{https://whyisyoung.github.io/BODMAS/}.

\bibitem[\citeproctext]{ref-yevsikov2026dllicious}
Yevsikov, Alexander, and Nir Nissim. 2026. {``{DLLicious}: Detection and
Explainability of Malicious DLL Files Using Novel Static Feature
Extraction Methods.''} \emph{Expert Systems with Applications} 320:
132188. \url{https://doi.org/10.1016/j.eswa.2026.132188}.

\end{CSLReferences}

\end{document}
